\documentclass[10pt,a4paper]{amsart}
\usepackage[margin=19mm]{geometry}
\usepackage{amsmath,amssymb,mathtools}
\usepackage[scr=rsfs,cal=euler]{mathalfa}
\usepackage{xfrac}
\usepackage[unicode,hidelinks,bookmarks=false]{hyperref}
\newcommand{\ie}{i.e.\ }
\newcommand{\cf}{cf.\ }
\newcommand{\resp}{respectively\ }
\newcommand{\Subsection}{\subsection}
\providecommand{\nobreakdash}{\nobreak\hbox{-}}
\setlength{\emergencystretch}{2em}
\multlinegap0pt
\usepackage{amssymb}
\usepackage{mathtools}
\usepackage{enumerate}
\newtheorem{lemma}{Lemma}
\newtheorem{theorem}{Theorem}
\newcommand{\C}{\mathbb{C}}
\newcommand{\N}{\mathbb{N}}
\newcommand{\R}{\mathbb{R}}
\title{Remling's Theorem for vector-valued discrete Schr\"odinger operators}
\author{Keshav Raj Acharya}
\date{Source: arXiv:2603.00299v1 (CC BY 4.0)}
\begin{document}
\maketitle

\noindent\emph{Source and licence.} Remling's Theorem for vector-valued discrete Schr\"odinger operators. Version arXiv:2603.00299v1 (CC BY 4.0), distributed by arXiv under the Creative Commons Attribution 4.0 International licence, http://creativecommons.org/licenses/by/4.0/, as recorded in the arXiv licence metadata for this exact version and shown on the abstract page https://arxiv.org/abs/2603.00299v1. Full licence notice: \texttt{licenses/arxiv-2603.00299-CC-BY-4.0.txt}.\par\smallskip

\noindent\textit{Editorial scope.} Counted unit: the article's theorem 4.8 (source0.tex lines 423--427). The article's complete original proof of it is delivered with it (source0.tex lines 568--618, one \texttt{proof} environment of 51 lines, which the article places away from the statement). No mathematics is written, repaired or re-derived: the statement and the proof are the article's own lines, in the article's own notation, and no retained line is substituted or re-flowed.  Retained uncounted as the source-local context the proof needs: lines 75--77; lines 348--350; lines 402--403; lines 462--464; lines 516--517; lines 558--559.  Nothing was omitted from any retained span, so the package carries no omission record.  No retained line contains a citation, so the wrapper carries no bibliography and no bibliography entry.  No retained line contains a figure environment, an image-inclusion command or drawing code, so no image is delivered and the package declares no asset dependency.  Editorial formatting only, in addition: an AMS \texttt{amsart} preamble that restates the article's own declarations and macros used by the retained lines, namely the environments \texttt{lemma}, \texttt{theorem} on separate counters, together with the standard packages those lines need.  The retained environments print in this document as Lemma 1, Theorem 1 in order of appearance, whereas the article prints the same items with the numbering of its own full document.  The complete native source of the article as distributed in the arXiv e-print for this exact version is bundled unchanged as \texttt{sources/195477/source0.tex}.
\begin{align}\label{ds}
y(n+1)+y(n-1)+B(n)y(n)=zy(n),\quad z\in \C 
\end{align}

\begin{equation} \label{hmp}
\omega_z(S) = \omega_{-\overline{z}}(-S). 
\end{equation} 

\begin{lemma} \label{lemma 4.6} Let $A \subset \R$ be a Borel set with $|A| < \infty $. Then  for any $c\in \mathbb B(0)$ 
$$ \lim_{y\rightarrow 0+} \sup_{M\in \mathbb H, \ \ S \subset \R} \Big| \int_A \omega_{c^*M(t+iy)c}(S) dt  - \int_A \omega_{c^*M(t)c}(S) dt \Big| = 0 .$$\end{lemma}

 \begin{equation} \label{eq 4.11}
 |\omega_{c^* M_1 c}(S)-\omega_{c^*M_2c}(S)|\leq \gamma(c^*M_1c,c^*M_2 c),
 \end{equation} for any $c \in \mathbb B(0)$ and any Borel set $S\subset\R.$ 

 Introduce $$P_{+}(n,z):= T_{+}(n,z)T_{\pm}(n-1,z) \dots T_{+}(1,z).$$ and $$ P_-(n,z) := T_{-}(n,z)T_{-}(n-1,z) \dots T_{-}(2,z).$$ Thus  \begin{equation} \label{eq 4.14}  \tilde  M_+ (n,z) =P_{+}(n,z) \tilde M_+(0,z), \    \ \tilde M_- (n,z) =P_{-}(n,z) \tilde M_-(1,z)    
 .\end{equation}

\begin{lemma} \label{lemma 4.11 } Let K be a compact subset of $\C^+.$ Then for any $W\in S_d,$
	 $$ \lim_{n\rightarrow \infty}	d_{\infty}(\tilde M_-(n, z),  P_-(n, z)W) = 0  $$  uniformly in $z\in K .$\end{lemma}

\begin{theorem} \label{theorem 4.8}  Let $\Sigma_{ac}^d$ denotes the essential support with full multiplicity of absolutely continuous part of the spectral measure of the half-line problem \eqref{ds}. Then for any $A \subset \Sigma_{ac}^d, \ |A| < \infty $ and $S \subset \R,$ we have 
\begin{equation*}
\lim_{N\rightarrow \infty}\Big ( \int_A \omega_{c^* \tilde M_-(N,t)c} (-S) dt -\int_A \omega_{c^* \tilde M_+(N,t)c}(S) dt\ \Big ) = 0,
\end{equation*} 
for all $c\in \mathbb B(0).$ Moreover, the convergence is uniform in $S .$\end{theorem}

\begin{proof}[Proof of Theorem \ref{theorem 4.8}] Let $ A \subset \Sigma_{ac}^d$ with $|A| < \infty$ and let $\epsilon >0$ be given. Partition   $A = A_0\cup A_1 \cup \dots \cup A_N$ of disjoint 
 subsets such that $|A_0| <\epsilon,  $ and $ A_j$ is bounded for $j\geq 1.$ We require that $M_+(t) = \lim_{y \rightarrow 0^+} M_+(t+iy)$ exists and $M_+(t) \in \mathcal S_d$ on $ \bigcup_{j=1}^{N} A_j .$ To find $A_j$ with these properties, first of all put $t \in A$ for which $M_+(t)$ does not exists or does not lie in $\mathcal S_d $ into $A_0.$ Clearly, $|A_0|\leq \epsilon.$ Then pick a sufficiently large compact subset $S \subset \mathcal S_d$ that contains the range of $M_+= M_+(0,t)$, and a compact subset $ t\in K \subset \R .$ Subdivide $S$ into finitely many disjoint subsets   $ S_1, S_2,  \dots ,S_n$  such that for all $j=1, \dots, n, $ we have that, for any $ Z_1, Z_2 \in S_j,$   $$ d_{\infty}(Z_1, Z_2) \leq \epsilon .$$ Then define $A_j = (A \setminus A_0)\cap \tilde M_+^{-1}(S_j)$ and $ M^j = \tilde M_+(t_j) $ for any fixed $ t_j \in A_j. $ Note that for each $j$, $\tilde M_+(t_j)$ is invertible. These $ M^j $ satisfy the  inequality:  \begin{equation} \label{eq 4.16} d_{\infty}(\tilde M_+(t), \ M^j) \leq \epsilon \end{equation} for all $t \in A_j,\quad j\geq 1 .$
 
 By Lemma \ref{lemma 4.6}, for any $c\in \mathbb B(0),$ there is a number $y>0$ such that, for arbitrary matrix-valued Herglotz function $M \in \mathbb H$, for any Borel subset $S\subset \R$ and for all $j =1,2, \dots, n$ we have the following estimate,
 
 \begin{equation} \label{eq 4.22}  \Big | \int_{A_j} \omega_{c^*  M (N,t+iy)c} (-S) dt - \int_{A_j} \omega_{c^*  M (N,t)c}(-S) dt\  \Big | \leq \epsilon |A_j| .\end{equation}
We can define $y$ for each value of $j$ and $c\in \mathbb B(0)$; so $y$ is a function of $j $ and $c$, $ y =y(j,c)$. However, by taking the infimum over the sets $ \{1,2, \dots, n\}$ and $\mathbb B(0)$, we can assume $y$ to be independent of $j$ and $c$.

Since $P_+(n,t) \in  \operatorname{Aut}(\mathcal S_d)$ we obtain from \eqref{eq 4.14} and \eqref{eq 4.16}
that
 \begin{equation} \label{eq 4.18} d_{\infty}(\tilde M_+(n,t), \ P_+(n,t) M^j) \leq \epsilon \end{equation}
 for all $t \in A_j,\quad  1\leq j \leq N .$

Using \eqref{eq 4.11} and integrating show that for arbitrary Borel set $S \subset \R$ and for any $c \in \mathbb B(0)$
  \begin{align*}  \Big | \int_{A_j} \omega_{c^* \tilde M_+(n,t)c} (S) \ dt -  & \int_{A_j} \omega_{c^*  P_+(n,t)M^jc}(S) \ dt\  \Big | \\   & \leq  \int_{A_j}  \Big | \omega_{c^* \tilde M_+(N,t)c} (S) - \omega_{c^*  P_+(n,t)M^jc}(S) \Big|\ dt \\ & \leq  \gamma ( c^* \tilde M_+(n,t)c, \ c^*P_+(n,t) M^jc ) \ |A_j| \\ & \leq d_{\infty}(\tilde M_+(n,t), \ P_+(n,t) M^j) \ |A_j|. \end{align*}
  Thus by using \eqref {eq 4.18}
   \begin{align} \label{ineq 1}  \Big | \int_{A_j} \omega_{c^* \tilde M_+(N,t)c} (S) \ dt -  & \int_{A_j} \omega_{c^*  P_+(n,t)M^jc}(S) \ dt \Big |  \leq \epsilon |A_j|. \end{align}
 Now using   \eqref{hmp}, we can obtain 
 \begin{equation} \label{omega} \omega_{c^*P_+(n,t)M^jc}(S) = \omega_{- \overline{c^*P_+(n,t)M^jc}}(-S).\end{equation}
 Let $W^j = T_+(1, t)(-M^j).$ Then $W^j \in \mathcal S_d .$  
By Lemma \ref{lemma 4.11 } there exists $n_0 \in \N $ such that
\begin{equation} \label{eq 4.26}  d_{\infty}(\tilde M_-(n,t+iy), \ P_-(n,t+iy)( T_+(1,z)(- M^j ) \leq \epsilon \end{equation}
 for all $ n\geq n_0$, $t \in A_j,\quad  1\leq j \leq N .$
Similarly, again by Lemma \ref{lemma 4.6} and on integration over $A_j$, we estimate
 \begin{equation} \label{eq 4.28} \Big | \int_{A_j} \omega_{c^* P_-(n, t+iy) W^jc}(-S) dt - \int_{A_j} \omega_{c^* \tilde M_-(n, t+iy)c}(-S) dt\  \Big | \leq \epsilon |A_j| ,\end{equation} for any $c \in \mathbb B(0), n\geq n_0.$
Now by triangle inequality we get,
\begin{align*}   &\Big | \int_{A_j} \omega_{c^* P_-(n, t) W^jc}(-S) dt -  \int_{A_j} \omega_{c^* \tilde M_-(n, t)c}(-S) dt\  \Big | 
\\ & \leq  \Big | \int_{A_j} \omega_{c^* P_-(n,t) W^j c}(-S) dt - \int_{A_j} \omega_{c^* \tilde P_-(n, t+iy)W^jc}(-S) dt \Big | \\  & +\Big |\int_{A_j} \omega_{c^* \tilde P_-(n, t+iy)W^jc}(-S) dt - \int_{A_j} \omega_{c^*  \tilde M_-(n, t+iy)c}(-S) dt  \Big | \\ & + \Big |
\int_{A_j} \omega_{c^*  \tilde M_-(n, t+iy)c}(-S) dt  -
\int_{A_j} \omega_{c^* \tilde M_-(n, t)c}(-S) dt\  \Big | .\end{align*} 
Again by  \eqref{eq 4.22}   and \eqref{eq 4.28}, for any $c \in \mathbb B(0),  n\geq n_0$
\begin{align} \label{ineq2}  &\Big | \int_{A_j} \omega_{c^* P_-(n,t) W^j c}(-S) dt -  \int_{A_j} \omega_{c^* \tilde M_-(n,t)c}(-S) dt\  \Big | \leq 3\epsilon |A_j| . 
\end{align} 
  Using $ P_-(n,t) W^j = - \overline {P_+(n,t)M^j}$, for all for all $c\in \mathbb B(0)$ if $n\geq n_0 $ we get
 
  \begin{align*} 
  &\Big | \int_{A_j} \omega_{c^* \tilde M_-(N,t)c} (-S) dt   -\int_{A_j} \omega_{c^* \tilde M_+(N,t)c}(S) dt\ \Big | \\ 
  & \leq 
  \Big | \int_{A_j} \omega_{c^* \tilde M_-(N,t)c} (-S) dt  - \int_{A_j} \omega_{c^* P_-(n,t) W^j c}(-S) dt \Big |\\ & + \Big | \int_{A_j} \omega_{-c^*  \overline {P_+(n,t)M^j }c }(-S) dt - \int_{A_j} \omega_{c^* \tilde M_+(N,t)c}(S) dt\ \Big |
  .\end{align*}  
 By using \eqref{omega} and   combining this with \eqref{ineq2}, \eqref{eq 4.26}, \eqref{eq 4.28} we have,
 
 \begin{align*} \Big | \int_{A_j} \omega_{c^* \tilde M_-(N,t)c} (-S) dt   -\int_{A_j} \omega_{c^* \tilde M_+(N,t)c}(S) dt\ \Big |  
   \leq 4\epsilon |A_j| + \epsilon
 . \end{align*}  
 Now taking the sum over $j=1,2, \dots,N$ and using $|A_0 |\leq \epsilon $ we get 
  
  \begin{align*}
  \Big | \int_{A} \omega_{c^* \tilde M_-(N,t)c} (-S) dt   -\int_{A} \omega_{c^* \tilde M_+(N,t)c}(S) dt\ \Big |  
   \leq 4\epsilon |A| + 2\epsilon.
  \end{align*}  \end{proof}

\end{document}
